\documentclass[a4paper,11pt]{jsarticle}
\usepackage{amsmath,amssymb,amsthm,mathtools}
\usepackage{geometry}
\geometry{margin=25mm}
\newtheorem{definition}{定義}[section]
\newtheorem{theorem}[definition]{定理}
\newtheorem{proposition}[definition]{命題}
\newtheorem{remark}[definition]{注意}
\newcommand{\Solver}{\mathcal{S}}
\newcommand{\Sel}{I}
\newcommand{\Ax}{S_{\mathrm{ax}}}
\newcommand{\Th}{\mathrm{Th}}
\newcommand{\Zset}{\mathcal{Z}}
\newcommand{\Pow}{\mathcal{P}}
\title{ECA_9\\$\mathcal{Z}$ の濃度追跡と連続体濃度の下界}
\author{}
\date{}
\begin{document}
\maketitle
\begin{abstract}
本稿では、ECA_6 および ECA_7 で確定した Solver、公理選択空間、理論写像、$\mathcal{Z}$ の定義を前提として、
$|\mathcal{Z}|\geq 2^{\aleph_0}$ をどのような条件の下で導出できるかを整理する。
これまでの検証では、$\mathcal{P}(\mathbb{N})\hookrightarrow\mathcal{P}(\mathcal{A}_0)$ は構成できる一方、
$\mathcal{P}(\mathbb{N})\hookrightarrow I$ や $\mathcal{P}(\mathbb{N})\hookrightarrow\mathcal{Z}$ は現行定義だけからは導出できないことを確認している。
そこで本稿では、$\mathcal{Z}$ 内に連続体濃度の異なる Solver を実際に構成するための条件を明示し、その条件から得られる濃度下界を定理として記述する。
CH の証明は本稿の目的としない。
\end{abstract}
\section{出発点となる ECA の構造}
ECA_6 により、公理候補全体を $\Ax$ とし、公理選択空間を
\[
\Sel=I\subseteq\Pow(\Ax)
\]
とする。Solver は
\[
\boxed{s=(\mathcal{A}_s,\mathcal{R}_s)}
\]
であり、$\mathcal{A}_s\in I$ を満たすものを ECA に適合する Solver とする。したがって、
\[
\Solver_I=\{s=(\mathcal{A}_s,\mathcal{R}_s)\mid\mathcal{A}_s\in I\}.
\]
また、$\mathfrak{A}:\Solver_I\to I$ を $\mathfrak{A}(s)=\mathcal{A}_s$ とする。
ECA_7 では
\[
\boxed{\mathcal{Z}=\{s\in\Solver_I\mid\mathcal{T}(s)\cong_{\mathrm{Th}}\mathrm{ZFC}\}}
\]
と定義した。したがって $\mathcal{Z}\subseteq\Solver_I$ である。
\section{これまでに得られている濃度関係}
可算な公理候補集合 $\mathcal{A}_0=\{a_1,a_2,a_3,\ldots\}$ を固定する。
任意の $A\subseteq\mathbb{N}$ に対して
\[
\mathcal{A}_A=\{a_n\in\mathcal{A}_0\mid n\in A\}
\]
を定義すると、
\[
\Gamma:\Pow(\mathbb{N})\to\Pow(\mathcal{A}_0),\qquad \Gamma(A)=\mathcal{A}_A
\]
は単射である。したがって
\[
|\Pow(\mathcal{A}_0)|=2^{\aleph_0}.
\]
しかし、$\Pow(\mathcal{A}_0)\subseteq\Pow(\Ax)$ や $I\subseteq\Pow(\Ax)$ だけから $\Pow(\mathbb{N})\hookrightarrow I$ は導かれない。さらに $\mathcal{Z}\subseteq\Solver_I$ であるため、$\Pow(\mathbb{N})\hookrightarrow\mathcal{Z}$ を得るには、$\mathcal{Z}$ 内で連続体濃度の族を実現する追加条件が必要である。
\section{$\mathfrak{A}|_{\mathcal{Z}}$ による濃度追跡}
\[
\boxed{\mathfrak{A}|_{\mathcal{Z}}:\mathcal{Z}\to I}
\]
を考え、その像を
\[
\operatorname{Im}(\mathfrak{A}|_{\mathcal{Z}})=\{\mathcal{A}_s\mid s\in\mathcal{Z}\}
\]
とする。任意の写像について
\[
|\operatorname{Im}(\mathfrak{A}|_{\mathcal{Z}})|\leq|\mathcal{Z}|
\]
が成立する。したがって、像の濃度が $2^{\aleph_0}$ 以上であることを示せれば、$|\mathcal{Z}|\geq2^{\aleph_0}$ が得られる。
ただし $\mathfrak{A}|_{\mathcal{Z}}$ は一般には単射ではない。$\mathcal{A}_s=\mathcal{A}_t$ でも $\mathcal{R}_s\neq\mathcal{R}_t$ ならば $s\neq t$ であり得るからである。
\section{連続体濃度の族を $\mathcal{Z}$ 内に構成する条件}
可算な $\mathcal{A}_0$ と固定された推論規則・構成制約 $\mathcal{R}_0$ を考える。任意の $A\subseteq\mathcal{A}_0$ に対して
\[
s_A=(A,\mathcal{R}_0)
\]
という Solver 候補を考える。
この候補が $\mathcal{Z}$ に属するには、少なくとも $A\in I$ および
\[
\mathcal{T}(s_A)\cong_{\mathrm{Th}}\mathrm{ZFC}
\]
が必要である。
\begin{definition}[$\mathcal{Z}$ 内連続体実現条件]
可算な $\mathcal{A}_0$ と固定された $\mathcal{R}_0$ に対して
\[
\boxed{\forall A\subseteq\mathcal{A}_0,\quad (A,\mathcal{R}_0)\in\mathcal{Z}}
\]
が成立するとき、これを $\mathcal{Z}$ 内連続体実現条件と呼ぶ。
\end{definition}
この条件は、単に $A\subseteq\mathcal{A}_0$ を構成できることではなく、すべての $A$ について ECA 適合性と ZFC との理論同型性を同時に要求する。
\section{連続体濃度下界の定理}
\begin{theorem}
$\mathcal{A}_0$ が可算無限集合であり、$\mathcal{Z}$ 内連続体実現条件が成立すると仮定する。このとき
\[
\boxed{|\mathcal{Z}|\geq2^{\aleph_0}}
\]
が成立する。
\end{theorem}
\begin{proof}
\[
\Phi:\Pow(\mathcal{A}_0)\to\mathcal{Z},\qquad \Phi(A)=(A,\mathcal{R}_0)
\]
と定義する。実現条件により $\Phi(A)\in\mathcal{Z}$ である。
$A\neq B$ ならば第一成分が異なるため
\[
(A,\mathcal{R}_0)\neq(B,\mathcal{R}_0),
\]
したがって $\Phi$ は単射である。可算無限集合 $\mathcal{A}_0$ に対して $|\Pow(\mathcal{A}_0)|=2^{\aleph_0}$ なので、
\[
2^{\aleph_0}=|\Pow(\mathcal{A}_0)|\leq|\mathcal{Z}|.
\]
\end{proof}
\section{この定理が示していること}
核心は、単なる
\[
\Pow(\mathbb{N})\hookrightarrow\Pow(\mathcal{A}_0)
\]
ではなく、
\[
\Pow(\mathcal{A}_0)\hookrightarrow\mathcal{Z}
\]
という新たな単射を構成することである。すなわち、
\[
\Pow(\mathbb{N})\cong\Pow(\mathcal{A}_0)\hookrightarrow\mathcal{Z}
\]
という連鎖を成立させる必要がある。
\section{現行 ECA 定義から自動的に導かれるか}
現行の定義からは
\[
\mathcal{Z}=\{s\in\Solver_I\mid\mathcal{T}(s)\cong_{\mathrm{Th}}\mathrm{ZFC}\}
\]
しか与えられていない。したがって
\[
\forall A\subseteq\mathcal{A}_0,\quad(A,\mathcal{R}_0)\in\mathcal{Z}
\]
は、この定義だけからは導出されない。
特に、
\[
\forall A\subseteq\mathcal{A}_0,\ A\subseteq\Ax
\]
と
\[
\forall A\subseteq\mathcal{A}_0,\ (A,\mathcal{R}_0)\in\mathcal{Z}
\]
は別の条件である。前者は公理候補としての集合論的構成を表すだけであり、後者は各選択から ZFC と理論同型な Solver が実現されることまで要求する。
したがって現段階で確定できるのは、
\[
\boxed{\text{$\mathcal{Z}$ 内連続体実現条件}\Longrightarrow |\mathcal{Z}|\geq2^{\aleph_0}}
\]
という条件付き結果である。
\section{より一般的な単射条件}
\begin{proposition}
$\mathcal{A}_0$ が可算無限であり、
\[
\Phi:\Pow(\mathcal{A}_0)\hookrightarrow\mathcal{Z}
\]
が単射であるならば、
\[
\boxed{|\mathcal{Z}|\geq2^{\aleph_0}}
\]
である。
\end{proposition}
\begin{proof}
$|\Pow(\mathcal{A}_0)|=2^{\aleph_0}$ であり、$\Phi$ が単射なので
\[
2^{\aleph_0}=|\Pow(\mathcal{A}_0)|\leq|\mathcal{Z}|.
\]
\end{proof}
\section{濃度追跡の現在位置}
\[
\mathcal{A}_0\cong\mathbb{N}
\]
から
\[
\boxed{|\Pow(\mathcal{A}_0)|=2^{\aleph_0}}
\]
までは通常の集合論として確定する。一方、
\[
\Pow(\mathcal{A}_0)\hookrightarrow I
\]
および
\[
\Pow(\mathcal{A}_0)\hookrightarrow\mathcal{Z}
\]
は現行 ECA 定義だけからは未確定である。
ただし、後者を構成できれば
\[
\boxed{|\mathcal{Z}|\geq2^{\aleph_0}}
\]
が直ちに得られる。
したがって現在の核心は、ECA の既存定義から $\mathcal{Z}$ 内の連続体濃度の Solver 族を構成できるか、という問題である。
\section{CH との切り分け}
本稿の結果
\[
|\mathcal{Z}|\geq2^{\aleph_0}
\]
は CH
\[
2^{\aleph_0}=\aleph_1
\]
を意味しない。また、そこから
\[
2^{\aleph_0}\leq\aleph_1
\]
を導くこともできない。
したがって、連続体濃度以上の $\mathcal{Z}$ の構成と CH の成立は別問題として扱う。
\section{次段階}
次に検証すべきなのは、
\[
s_A=(A,\mathcal{R}_A)
\]
をどのように定めれば
\[
\mathcal{A}_{s_A}=A,
\qquad
\mathcal{T}(s_A)\cong_{\mathrm{Th}}\mathrm{ZFC}
\]
を同時に保証できるかである。
ここでは、単に公理選択を増やすのではなく、
\[
\mathcal{Z}=\mathcal{T}^{-1}([\mathrm{ZFC}]_{\mathrm{Th}})
\]
という $\mathcal{Z}$ の構造そのものに対して、連続体濃度の族を実現できるかを調べる。
\end{document}
